\chapter{Summary}
\label{ch:sum}

A theoretical framework was established to model $N$ dielectric trapped particles in an optical cavity.
This model provides a robust methodology for determining the equilibrium positions 
    and characterizing the stability of those positions, specifically for the single particle case.

Also notable we discovered that with a single particle there are configurations with multiple stable equilibria. 
When examining the detuning vs the $z$ displacement, 
    we get a continuous curve, see Figure \ref{fig:sp:stability} and \ref{fig:sp:kappas}.
Contrary to initial intuition,
    all equilibrium positions along this curve were found to be stable.
This result challenges the expected bistable behavior,
    which only emerges when applying an adiabatic approximation (see Appendix \ref{apdx:sp:fcf}).
This approximation fails, 
    since the cavity needs sufficient time to reach its steady state.
For $2$ particles we found that for a fixed detuning the solutions are found as discrete points along a curve in configuration space. 
These results are also novel 
    since it is typically assumed that $\kappa$ is very large.
Since we also consider smaller values of $\kappa$ our predictions are very relevant for the field for further research,
    which is getting towards better cavities, 
    where this domain can be explored.

This thesis establishes the fundamental principles required to transition this system into a quantum mechanical framework and explore emergent effects.
Further research options are discussed extensively in Chapter \ref{ch:otl}.